%% LyX 2.3.6.1 created this file.  For more info, see http://www.lyx.org/.
%% Do not edit unless you really know what you are doing.
\documentclass[11pt,a4paper,english]{article}
\usepackage{mathpazo}
\renewcommand{\sfdefault}{cmss}
\renewcommand{\ttdefault}{cmtt}
\usepackage[T1]{fontenc}
\usepackage[utf8]{inputenc}
\usepackage{xcolor}
\definecolor{shadecolor}{rgb}{1, 0, 0}
\usepackage{babel}
\usepackage{float}
\usepackage{framed}
\usepackage{url}
\usepackage{graphicx}
\usepackage[unicode=true,pdfusetitle,
 bookmarks=true,bookmarksnumbered=false,bookmarksopen=false,
 breaklinks=false,pdfborder={0 0 0},pdfborderstyle={},backref=page,colorlinks=false]
 {hyperref}

\makeatletter

%%%%%%%%%%%%%%%%%%%%%%%%%%%%%% LyX specific LaTeX commands.
\special{papersize=\the\paperwidth,\the\paperheight}

%% Because html converters don't know tabularnewline
\providecommand{\tabularnewline}{\\}

%%%%%%%%%%%%%%%%%%%%%%%%%%%%%% User specified LaTeX commands.


%------------------------------------------------------------------------------
%	REQUIRED PACKAGES AND  CONFIGURATIONS
%------------------------------------------------------------------------------
% PACKAGES FOR TITLES
\usepackage{titlesec}
\usepackage{color}
% PACKAGES FOR LANGUAGE AND FONT
%\usepackage[utf8]{inputenc}
%\usepackage[english]{babel}

% PACKAGES FOR IMAGES
\usepackage{graphicx}
\graphicspath{{Images/}}
\usepackage{eso-pic} % For the background picture on the title page
\usepackage{subfig} % Numbered and caption subfigures using \subfloat
\usepackage{caption} % Coloured captions
\usepackage{transparent}

% STANDARD MATH PACKAGES
\usepackage{amsmath}
\usepackage{amsthm}
\usepackage{bm}
\usepackage[overload]{empheq}  % For braced-style systems of equations

% PACKAGES FOR TABLES
\usepackage{tabularx}
\usepackage{longtable} % tables that can span several pages
\usepackage{colortbl}

% PACKAGES FOR ALGORITHMS (PSEUDO-CODE)
\usepackage{algorithm}
\usepackage{algorithmic}

% PACKAGES FOR REFERENCES & BIBLIOGRAPHY
%\usepackage[colorlinks=true,linkcolor=black,anchorcolor=black,citecolor=black,filecolor=black,menucolor=black,runcolor=black,urlcolor=black]{hyperref} % Adds clickable links at references
\usepackage{cleveref}
\usepackage[square, numbers, sort&compress]{natbib} % Square brackets, citing references with numbers, citations sorted by appearance in the text and compressed
%\bibliographystyle{plain} % You may use a different style adapted to your field

% PACKAGES FOR THE APPENDIX
\usepackage{appendix}

% PACKAGES FOR ITEMIZE & ENUMERATES 
\usepackage{enumitem}

% OTHER PACKAGES
\usepackage{amsthm,thmtools,xcolor} % Coloured "Theorem"
\usepackage{comment} % Comment part of code
\usepackage{fancyhdr} % Fancy headers and footers
\usepackage{lipsum} % Insert dummy text
\usepackage{tcolorbox} % Create coloured boxes (e.g. the one for the key-words)

% Configuration package
\usepackage[bottom=2.0cm,top=2.0cm,left=2.0cm,right=2.0cm]{geometry}
\raggedbottom 

% Create color blue (-> manuale grafica coordinata:  https://www.polimi.it/fileadmin/user_upload/il_Politecnico/grafica-coordinata/2015_05_11_46xy_manuale_grafica_coordinata.pdf)
\definecolor{blue}{cmyk}{0.4,0.1,0,0.4}
%\definecolor{shadecolor}{cmyk}{0.4,0.1,0,0.4}
\usepackage{xcolor}
\colorlet{shadecolor}{blue!60}
\colorlet{cyan}{blue!40}
% Custom theorem environments
\declaretheoremstyle[
  headfont=\color{blue}\normalfont\bfseries,
  bodyfont=\color{black}\normalfont\itshape,
]{colored}

\captionsetup[figure]{labelfont={color=blue}} % Set colour of the captions
\captionsetup[table]{labelfont={color=blue}} % Set colour of the captions
\captionsetup[algorithm]{labelfont={color=blue}} % Set colour of the captions

\theoremstyle{colored}
\newtheorem{theorem}{Theorem}[section]
\newtheorem{proposition}{Proposition}[section]

% Enhances the features of the standard "table" and "tabular" environments.
\newcommand\T{\rule{0pt}{2.6ex}}
\newcommand\B{\rule[-1.2ex]{0pt}{0pt}}

% Algorithm description
\newcounter{algsubstate}
\renewcommand{\thealgsubstate}{\alph{algsubstate}}
\newenvironment{algsubstates}{
    \setcounter{algsubstate}{0}%
    \renewcommand{\STATE}{%
    \stepcounter{algsubstate}%
    \Statex {\small\thealgsubstate:}\space}
    }{}
    
% Custom theorem environment
\newcolumntype{L}[1]{>{\raggedright\let\newline\\\arraybackslash\hspace{0pt}}m{#1}}
\newcolumntype{C}[1]{>{\centering\let\newline\\\arraybackslash\hspace{0pt}}m{#1}}
\newcolumntype{R}[1]{>{\raggedleft\let\newline\\\arraybackslash\hspace{0pt}}m{#1}}

% Custom itemize environment
\setlist[itemize,1]{label=$\bullet$}
\setlist[itemize,2]{label=$\circ$}
\setlist[itemize,3]{label=$-$}
\setlist{nosep}

% Create command for background pic
\newcommand\BackgroundPic{% Adding background picture
	\put(237,365){
	    \parbox[b][\paperheight]{\paperwidth}{%
	    \vfill
		\centering
		\transparent{0.4}
		\includegraphics[width=0.44\paperwidth]{raggiera_polimi}%
		\vfill}
		}
}

% Set indentation
\setlength\parindent{0pt}

% Custom title commands
\titleformat{\section}
{\color{blue}\normalfont\Large\bfseries}
{\color{blue}\thesection.}{1em}{}
\titlespacing*{\section}
{0pt}{3.3ex}{3.3ex}

\titleformat{\subsection}
{\color{blue}\normalfont\large\bfseries}
{\color{blue}\thesubsection.}{1em}{}
\titlespacing*{\subsection}
{0pt}{3.3ex}{3.3ex}


\titleformat{\subsubsection}
{\color{blue}\normalfont\large\bfseries}
{\color{blue}\thesubsubsection.}{1em}{}
\titlespacing*{\subsubsection}
{0pt}{3.3ex}{3.3ex}


% Custom headers and footers
\pagestyle{fancy}
\fancyhf{}
      
\fancyfoot{}
\fancyfoot[C]{\thepage} % page
\renewcommand{\headrulewidth}{0mm} % headrule width
\renewcommand{\footrulewidth}{0mm} % footrule width


\usepackage{anyfontsize}

\usepackage{bold-extra}
%\usepackage{fix-cm}
%\usepackage{cm-super}
%%
%\makeatletter
%%

\makeatother

\begin{document}
\AddToShipoutPicture*{\BackgroundPic}

\hspace{-0.6cm}\includegraphics[width=0.6\textwidth]{Images/logo_polimi_ing_indinf}

\vspace{-1mm}

\textcolor{blue}{\LARGE{}Title}{\LARGE\par}

\vspace{0.5cm}

\textbf{\textsc{\textcolor{blue}{\small{}Tesi di Laurea Magistrale
in }}}{\small\par}

\textbf{\textsc{\textcolor{blue}{\small{}Electronics Engineering -
Ingegneria Elettronica}}}{\small\par}

\vspace{0.3cm}

\textbf{\large{}Nome Cognome, codice persona}{\large\par}

\vspace{11pt}

\rule{1\columnwidth}{0.4pt}

\noindent\begin{minipage}[c][1\totalheight][t]{1\columnwidth}%
\medskip{}

\enskip{}%
\begin{minipage}[c]{0.24\textwidth}%
\begin{center}
\begin{minipage}[t][1\totalheight][c]{0.9\textwidth}%
\textbf{\scriptsize{}Advisor:}{\scriptsize\par}

{\scriptsize{}Prof. Nicola Lusardi}{\scriptsize\par}

\medskip{}

\textbf{\scriptsize{}Co-Advisors:}{\scriptsize\par}

{\scriptsize{}Dr. Nicola Corna}{\scriptsize\par}

{\scriptsize{}Dr. Fabio Garzetti}{\scriptsize\par}

\medskip{}

\textbf{\scriptsize{}Academic Year:}{\scriptsize\par}

{\scriptsize{}2021-2022 }{\scriptsize\par}%
\end{minipage}
\par\end{center}%
\end{minipage}\enskip{}\enskip{}%
\begin{minipage}[c]{0.74\columnwidth}%
\textbf{\textcolor{blue}{\footnotesize{}Abstract: }}\textcolor{black}{\footnotesize{}Here
goes the Abstract in English of your thesis (in article format) followed
by a list of keywords. The Abstract is a concise summary of the content
of the thesis (single page of text) and a guide to the most important
contributions included in your thesis. The Abstract is the very last
thing you write. It should be a self-contained text and should be
clear to someone who hasn’t (yet) read the whole manuscript. The Abstract
should contain the answers to the main research questions that have
been addressed in your thesis. It needs to summarize the motivations
and the adopted approach as well as the findings of your work and
their relevance and impact. The Abstract is the part appearing in
the record of your thesis inside POLITesi, the Digital Archive of
PhD and Master Theses (Laurea Magistrale) of Politecnico di Milano.
The Abstract will be followed by a list of four to six keywords. Keywords
are a tool to help indexers and search engines to find relevant documents.
To be relevant and effective, keywords must be chosen carefully. They
should represent the content of your work and be specific to your
field or sub-field. Keywords may be a single word or two to four words.}{\footnotesize\par}

\textcolor{black}{\footnotesize{}some citations to check bibliography
\cite{AXI_ref_guide}}{\footnotesize\par}%
\end{minipage}%
\end{minipage}

\vspace{15pt}

\noindent\begin{minipage}[t]{1\columnwidth}%
\begin{shaded}%
\textcolor{white}{\enskip{}}\textbf{\textcolor{white}{Key-words:}}\textcolor{white}{{}
insert keword list }\end{shaded}%
\end{minipage}

\vspace{12pt}


\section{Introduction}



\section{State of the Art}



\section{Main Work}



\section{Measures}



\section{Conclusions}


\bibliographystyle{plain}
\bibliography{bibtex}

\input{\string"A.AppendixA/Appendix A.tex\string"}

\clearpage{}

\input{\string"*Abstract ita/*Abstract ita.tex\string"}

\noindent\begin{minipage}[t]{1\columnwidth}%
\begin{shaded}%
\textcolor{white}{\enskip{}}\textbf{\textcolor{white}{Parole chiave:}}\textcolor{white}{{}
inserisci lista di parole chiave}\end{shaded}%
\end{minipage}


\section*{Acknowledgements}


\clearpage{}

\section*{Guide to Lyx colors}

This guide addresses the problem of how to get in $LtX$ the same
graphics effects present in the Polimi Thesis template, which can
be found here \url{https://it.overleaf.com/latex/templates/article-format-thesis-scuola-di-ingegneria-industriale-e-dellinformazione-politecnico-di-milano/vtqgsrqwzdmy}

\subsection{Tables}

Very fancy tables can be created, but \LaTeX code must be embedded
to add background colors, moreover, to embed mathemathical formulas
and greek letters, also latex code must be embedded

\noindent\colorbox{red}{\begin{minipage}[t]{1\columnwidth}%
\textcolor{white}{Pay attention: here the floatboxes are forces to
a fixed position, tick on ``use default placement'' in float settings
if u copy/paste this tables}%
\end{minipage}}

\begin{table}
\begin{centering}
\begin{tabular}{|c|c|c|c|}
\hline 
\rowcolor{blue!40} & \textbf{column1} & \textbf{column2} & \textbf{column3\T\B}\tabularnewline
\hline 
\hline 
\textbf{row1} & 1 & 2 & 3\tabularnewline
\hline 
\textbf{row2} & $\alpha$ & $\beta$ & $\mathrm{\gamma}$\tabularnewline
\hline 
\textbf{row3} & alpha & beta & gamma\tabularnewline
\hline 
\end{tabular}\caption{}
\par\end{centering}
\end{table}

\begin{center}
\begin{table}[H]
\centering{}%
\begin{tabular}{|c|c|c|c|c|c|}
\hline 
 \rowcolor{blue!40} & \textbf{column1} & \textbf{column2} & \textbf{column3} & \textbf{column4} & \textbf{column5\T\B}\tabularnewline
\hline 
\hline 
\textbf{Row1} & 1 & 2 & 3 & 4 & 5\textbf{\T\B}\tabularnewline
\textbf{Row2} & $\alpha$ & $\alpha$ & $\alpha$ & $\alpha$ & $\alpha$\textbf{\T\B}\tabularnewline
\textbf{Row3} & alpha & alpha & alpha & alpha & alpha\textbf{\T\B}\tabularnewline
\hline 
\end{tabular}\caption{Tables with Highlighted columns}
\end{table}
\par\end{center}

\subsection{Algorithms}

\begin{algorithm}[H] 
\label{alg:example} 
\caption{Name of the Algorithm} \label{alg:var} \label{protocol1} 
\begin{algorithmic}[1] 
\STATE Initial instructions 
\FOR{$for-condition$} 
\STATE{Some instructions} 
\IF{$if-condition$} 
\STATE{Some other instructions} 
\ENDIF 
\ENDFOR 
\WHILE{$while-condition$} 
\STATE{Some further instructions} 
\ENDWHILE 
\STATE Final instructions 
\end{algorithmic} 
\end{algorithm} 
\end{document}
